\documentclass[./main]{subfiles}

\begin{document}
  \chapter{Complexité de comptage.}

  \section{Problèmes de comptage.}

  \begin{exm}
    Les problèmes \pbSharpSat,\footnote{Le problème \pbSharpSatD\ vu au chapitre précédent est la version \textit{décision} du problème \pbSharpSat\ vu ici.} \pbSharpCliques, \pbSharpCouplagesParfaits\ sont des problèmes de comptage.
    \showSharpSat
    \showSharpCliques
    \showSharpCouplagesParfaits
  \end{exm}

  \begin{defn}
    On définit $\intro*\sharpp$ comme la classes des fonctions $f : \Sigma^\star \to \mathds{N}$ telles qu'il existe $A \in \p$ et un polynôme $p$ tel que 
     \[
       f(x) = \# \{y \in \Sigma^{\le p(n)}  \mid \code{x,y} \in A\} 
    .\]
  \end{defn}

  Cette définition est analogue à $\np$ pour les problèmes de comptage.

  \begin{defn}
    Le \textit{\textbf{permanent}} d'une matrice $A$ de taille $n \times n$ est défini comme
    \[
      \intro*\perm(A) := \sum_{\sigma \in \mathfrak{S}_n} a_{1, \sigma(1)} a_{2,\sigma(2)} \cdots a_{n,\sigma(n)}
    .\]
    (Il correspond au déterminant sans le facteur $\epsilon(\sigma)$, signature de la permutation $\sigma$.)
  \end{defn}

  \begin{exm}
    Pour des matrices à entrée dans $\mathds{N}$, on a ${\perm} \in \sharpp$ car 
    \[
      \perm(A) = \# \{ (\sigma, y)  \mid \underbrace{1 \le y \le a_{1, \sigma(1)} a_{2,\sigma(2)} \cdots a_{n,\sigma(n)}}_{\text{se vérifie en temps polynomial}} \} 
    .\]

    On peut interpréter le permanent en combinatoire : pour une matrice~$A \in \{\texttt{0},\texttt{1}\}^{n^2}$, le permanent $\perm(A)$ compte le nombre de couplages parfaits dans le graphe biparti correspondant ($A$ est la matrice d'adjacence de $G$).
    En effet, à une permutation $\sigma \in \mathfrak{S}_n$ correspond à un couplage parfait comme montré ci-dessous.


    \begin{figure}[H]
      \centering
      \begin{tikzpicture}[node distance=0.5cm]
        \node               (a1) {$\bullet$};
        \node[below= of a1] (a2) {$\bullet$};
        \node[below= of a2] (a3) {$\bullet$};
        \node[below= of a3] (a4) {$\bullet$};
        \node[right=1cm of a1] (b1) {$\bullet$};
        \node[below= of b1] (b2) {$\bullet$};
        \node[below= of b2] (b3) {$\bullet$};
        \node[below= of b3] (b4) {$\bullet$};
        \draw[thick, nicered] (a1) to (b1);
        \draw (a1) to (b3);
        \draw (a2) to (b1);
        \draw[thick, nicered] (a2) to (b3);
        \draw[thick, nicered] (a3) to (b2);
        \draw (a3) to (b4);
        \draw (a4) to (b2);
        \draw[thick, nicered] (a4) to (b4);
      \end{tikzpicture}
      \quad\quad
      \quad\quad
      \begin{tikzpicture}[node distance=1cm]
        \node               (a1) {$\bullet$};
        \node[right=of a1]  (a2) {$\bullet$};
        \node[below=of a1]  (a4) {$\bullet$};
        \node[below=of a2]  (a3) {$\bullet$};
        \draw[->, thick, nicered] (a1) to[loop left] (a1);
        \draw[->, thick, nicered] (a4) to[loop left] (a4);
        \draw[->, thick, nicered] (a2) to[bend right=20] (a3);
        \draw[->, thick, nicered] (a3) to[bend right=20] (a2);
        \draw[->] (a2) to (a1);
        \draw[->] (a1) to (a3);
        \draw[->] (a3) to (a4);
        \draw[->] (a4) to (a2);
      \end{tikzpicture}

      $\sigma : \big( 1 \mapsto 1 \quad 2 \mapsto 3 \quad 3 \mapsto 2 \quad 4 \mapsto 4 \big) $
      \caption{Interprétation combinatoire du permanent}
    \end{figure}

    Le permanent compte aussi le nombre de couvertures en cycle dans les graphes orientés.
    Les cycles de la permutation $\sigma \in \mathfrak{S}_n$ correspondent exactement aux cycles du graphe.
  \end{exm}

  \begin{rmk}
    \begin{itemize}
      \item La classe $\sharpp$ est "plus dure" que $\np$ : on a $\np \subseteq \oracle{\p}{\sharpp}$.
      \item La classe $\sharpp$ est "plus simple" que $\pspace$ : on a $\sharpp \subseteq \fpspace$ où $\intro*\fpspace$ est l'ensemble des \textit{fonctions} calculables en temps polynomial.
      \item On a même que la classe $\sharpp$ est plus dure que $\ph$ (théorème de Toda).
    \end{itemize}
  \end{rmk}

  \section{La $\sharpp$-complétude.}

  \begin{defn}
    On dit qu'une fonction $f \in \sharpp$ est \textit{\textbf{$\sharpp$-complète}} si, pour tout $g \in \sharpp$, il existe des fonctions $I$ et $O$ calculables en temps polynomial telles que
    \[
    g(x) = O(f(I(x)))
    .\]
    Dans ce cas, je noterai $(I, O) : g \intro*\redS f$.

    Si $O$ est la fonction identité alors on dit que la réduction est \textit{\textbf{parcimonieuse}}.
  \end{defn}

  La plupart des réductions entre problèmes (de comptages associés à des problèmes) $\np$-complets sont parcimonieuses.
  Ainsi, \pbSharpSat, \pbSharpChemHam, \textit{etc} sont $\sharpp$-complets par des réductions parcimonieuses.

  \showSharpChemHam


  On peut composer des réductions.

  \begin{lem}
    Si $(I_1, O_1) : f_2 \redS f_1$ et $(I_2, O_2) : f_3 \redS f_2$ alors on a
    \[
    O_2(O_1(f_1(I_1(I_2(x))))) = O_2(f_2(I_2(x))) = f_3(x)
    ,\] 
    d'où $(I_1 \circ I_2, O_2 \circ O_1) : f_3 \redS f_1$.
  \end{lem}

  \begin{thm}
    Le problème $\pbSharpCircuitSat$ est  $\sharpp$-complet.
    \showSharpCircuitSat
  \end{thm}
  \begin{prv}
    Soit $f \in \sharpp$.
    Il existe $A \in \p$ et $p$ tels que \[
      f(x) = \# \{y \in \{\texttt{0},\texttt{1}\}^{p(n)}  \mid \code{x, y} \in A\} 
    ,\]
    que l'on peut réécrire en
    \[
      f(x) = \# \{y \in \{\texttt{0},\texttt{1}\}^{p(n)}  \mid C_n(x,y) = \texttt{1}\} 
    ,\] 
    où $C_n$ est un circuit à $n + p(n)$ variables, car $A \in \p$. Or, c'est
    une instance de $\pbSharpCircuitSat$.
  \end{prv}

  \begin{thm}
    Le problème $\pbSharpSat$ est $\sharpp$-complet.
  \end{thm}
  \begin{prv}
    On réalise une réduction parcimonieuse à partir de $\pbSharpCircuitSat$.
    On introduit une variable booléenne  $z_g$ à chaque porte $g$ du circuit.
    Ainsi, pour une porte $\wedge$ sur $g_1$ et $g_2$, on peut exprimer "$z_g = z_{g_1} \wedge z_{g_2}$" avec la formule
    \[
      (z_g \vee \bar{z}_{g_1} \vee \bar{z}_{g_2}) \wedge (\bar{z}_g \vee z_{g_1}) \wedge (\bar{z}_g \vee z_{g_2})
    .\]
    (C'est la même réduction que $\pbCircuitSat \redP \pbSat$.)
    La réduction est parcimonieuse (une unique valeur possible pour chaque~$z_g$).
  \end{prv}

  \subsubsection{Est-ce que $\pbSharpCouplagesParfaits$ est $\sharpp$-complet ?}

  On sait que $\pbSharpCouplagesParfaits$ ne peut pas être  $\sharpp$-complet pour les réductions parcimonieuses (si $\p \neq \np$) car sinon on aurait
  \[
  \pbSharpSat(\phi) = \pbSharpCouplagesParfaits(I(\phi))
  ,\]
  or on peut décider si $\pbSharpCouplagesParfaits(I(\phi)) \ge 1$ et donc on pourrait décider $\pbSat$ en temps polynomial.

  \begin{thm}[Valiant]
    Compter le nombre de couplages parfaits dans un graphe biparti est $\sharpp$-complet.
  \end{thm}

  Pour montrer ce théorème, on montre d'abord que $\perm$ est $\sharpp$-complet.

  Le permanent, en toute généralité, s'interprète comme
  \begin{enumerate}
    \item la somme des poids des couplages parfaits (défini comme le produit des poids des arêtes) dans un graphe biparti ;
    \item la somme des poids des couvertures par cycles dans un graphe orienté.
  \end{enumerate}

  Le \textit{\textbf{gadget du XOR}} est le gadget ci-dessous.

  \begin{figure}[H]
    \centering
    \begin{tikzpicture}[node distance=2cm]
      \node (a) {$\bullet$} ;
      \node[above right= of a] (b0) {} ;
      \node[left=0.1cm of b0, inner sep=0pt] (bl) {} ;
      \node[right=0.1cm of b0, inner sep=0pt] (br) {} ;
      \draw[line width=5pt, line cap=round] (bl.center) -- (br.center);
      \node[inner sep=7pt, fit=(br)(bl)] (b) {};
      \node[inner sep=7pt, fit=(bl)] (bl) {};
      \node[inner sep=7pt, fit=(br)] (br) {};
      \node[below right= of a] (c0) {} ;
      \node[left=0.1cm of c0, inner sep=0pt] (cl) {} ;
      \node[right=0.1cm of c0, inner sep=0pt] (cr) {} ;
      \draw[line width=5pt, line cap=round] (cl.center) -- (cr.center);
      \node[inner sep=7pt, fit=(cr)(cl)] (c) {};
      \node[inner sep=7pt, fit=(cl)] (cl) {};
      \node[inner sep=7pt, fit=(cr)] (cr) {};
      \node[below right= of b0] (d) {$\bullet$} ;
      \node[above  left= of a] (u1) {$\bullet$} ;
      \node[above right= of d] (v1) {$\bullet$} ;
      \node[below  left= of a] (v2) {$\bullet$} ;
      \node[below right= of d] (u2) {$\bullet$} ;
      \node[below=0cm of a] {$a$};
      \node[above right=-0.25cm and -0.25cm of b] {$b$};
      \node[below right=-0.25cm and -0.25cm of c] {$c$};
      \node[below=0cm of d] {$d$};
      \node[below=0cm of u1] {$1$};
      \node[below=0cm of v1] {$1'$};
      \node[below=0cm of u2] {$2$};
      \node[below=0cm of v2] {$2'$};
      \draw[->] (u1) to (a);
      \draw[->] (a) to (v2);
      \draw[->] (d) to (v1);
      \draw[->] (u2) to (d);
      \draw[->] (a) to[bend left=20] (bl);
      \draw[->] (bl) to[bend left=20] (a);
      \draw[->] (d) to[bend left=20] (br);
      \draw[->] (br) to[bend left=20] (d);
      \draw[->] (d) to[bend left=20] node[midway,fill=white,text=nicered] {$3$} (cr);
      \draw[->] (cr) to[bend left=20] node[midway,fill=white,text=nicered] {$2$} (d);
      \draw[->] (a) to[bend right=20] node[midway,fill=white,text=nicered] {$-1$} (cl);
      \draw[->] (b) to[loop above] node[above,nicered] {$-1$} (b);
      \draw[->] (c) to[loop below] (c);
      \draw[->] (b) to[bend left=15] node[midway, fill=white, inner sep=3pt] {} (c);
      \draw[->] (c) to[bend left=15] node[midway, fill=white, inner sep=3pt] {} (b);
      \draw[->] (a) to node[near start,fill=white,text=nicered] {$-1$} (d);
    \end{tikzpicture}
    \caption{Gadget du XOR}
  \end{figure}

  On s'autorisera à brancher le gadget du XOR comme ci-dessous, en retirant les arêtes $1 1'$ et  $2 2'$.
  On note $G'$ le graphe ainsi obtenu.

  \begin{adjustbox}{center}
    \begin{tikzpicture}
      \node               (a1) {$\bullet$};
      \node[right=of a1]  (a2) {$\bullet$};
      \node[below=of a1]  (a4) {$\bullet$};
      \node[below=of a2]  (a3) {$\bullet$};
      \draw[->] (a1) to node[midway, inner sep=0pt] (m1) {} (a2);
      \draw[->] (a4) to node[midway, inner sep=0pt] (m2) {} (a3);
      \node[above=0cm of a1] {$1$};
      \node[above=0cm of a2] {$1'$};
      \node[below=0cm of a3] {$2'$};
      \node[below=0cm of a4] {$2$};
      \draw[nicered, thick] (m1) to node[midway, fill=white, text=nicered, inner sep=0pt] {$\bigoplus$} (m2);
    \end{tikzpicture}
  \end{adjustbox}

  On a que \[
    \sum_{\substack{\text{$C$ couverture par}\\ \text{cycles de $G'$}}} \mathrm{poids}(C) = 
    4 \times \sum_{\substack{
      \text{$C$ couverture par}\\
    \text{cycles de $G$ qui}\\
      \text{contient $11'$ \textit{\textbf{ou}}}\\
      \text{\textit{\textbf{exclusif}} $22'$}
    }} \mathrm{poids}(C)
  .\]
  On l'admet, cela pourra être vu en TD.

  Le \textit{\textbf{gadget de clause}} (noté $G_c$) est le gadget en figure~\ref{fig:gadget-clause}, représenté ici pour la clause $x \vee y \vee \bar{z}$.

  \begin{figure}
    \centering
    \begin{tikzpicture}[node distance=3cm]
      \node (base) {\phantom{$\bullet$}};
      \node[above=1.3cm of base] (A) {$\bullet$};
      \node[below left =of base] (B) {$\bullet$};
      \node[below right=of base] (C) {$\bullet$};
      \node[below=0.3cm of base] (D) {$\bullet$};
      \path (B) to[bend right=30] node[midway] (E) {$\bullet$} (C);
      \draw[->] (A) to[bend right=50] node[midway, inner sep=0pt] (AB) {} (B);
      \draw[->] (C) to[bend right=50] node[midway, inner sep=0pt] (CA) {} (A);
      \draw[->] (B) to[bend right=50] node[midway, inner sep=0pt] (BC) {} (C);
      \draw[->] (C) to[bend left=5] (B);
      \draw[->] (B) to[bend right=10] (E);
      \draw[->] (E) to[bend right=10] (C);
      \draw[->] (E) to[loop above] (E);
      \draw[->] (B) to[bend right=20] (D);
      \draw[->] (D) to[bend right=20] (B);
      \draw[->] (A) to[bend right=20] (D);
      \draw[->] (D) to[bend right=20] (A);
      \draw[->] (C) to[bend right=20] (D);
      \draw[->] (D) to[bend right=20] (C);
      \node[above left=1.5cm and 1.5cm of AB] (W0) {};
      \node[above right=1.5cm and 1.5cm of CA] (V0) {};
      \node[below=1.5cm and 1.5cm of BC] (U0) {};
      \node[above right=0.25cm and 0.25cm of W0] (W1) {$\bullet$};
      \node[below left=0.25cm and 0.25cm of W0] (W2) {$\bullet$};
      \node[above left=0.25cm and 0.25cm of V0] (V1) {$\bullet$};
      \node[below right=0.25cm and 0.25cm of V0] (V2) {$\bullet$};
      \node[right=0.6cm of U0] (U1) {$\bullet$};
      \node[left=0.6cm of U0] (U2) {$\bullet$};
      \draw[->] (U1) to[bend left]  node[very near end, below left] {\small $z = \texttt{1}$} (U2);
      \draw[->] (U1) to[bend right] node[midway, inner sep=0pt] (U) {} node[very near end, above left] {\small $z = \texttt{0}$} (U2);
      \draw[->] (U2) to (U1);
      \draw[->] (V1) to[bend left] node[very near end, right=0.2cm] {\small $y = \texttt{0}$} (V2);
      \draw[->] (V1) to[bend right] node[midway, inner sep=0pt] (V) {} node[very near end, below=0.2cm] {\small $y = \texttt{1}$} (V2);
      \draw[->] (V2) to (V1);
      \draw[->] (W1) to[bend left] node[midway, inner sep=0pt] (W) {} node[very near end, below=0.2cm] {\small $x = \texttt{1}$} (W2);
      \draw[->] (W1) to[bend right] node[very near end, left=0.2cm] {\small $x = \texttt{0}$} (W2);
      \draw[->] (W2) to (W1);
      \draw (U) to node[fill=white,text=black,midway,inner sep=0pt, sloped] (UP) {$\bigoplus$} (BC);
      \draw (V) to node[fill=white,text=black,midway,inner sep=0pt, sloped] (UP) {$\bigoplus$} (CA);
      \draw (W) to node[fill=white,text=black,midway,inner sep=0pt, sloped] (UP) {$\bigoplus$} (AB);
    \end{tikzpicture}
    \caption{Gadget de la clause $x \vee y \vee \bar{z}$}
    \label{fig:gadget-clause}
  \end{figure}

  La propriété du gadget de clause est la suivante :
  pour tout sous-ensemble strict de $X$ de l'ensemble $Y$ des trois arcs extérieurs, il existe une couverture par cycles unique qui contient uniquement les arcs de $X$ parmi ceux de $Y$. 
  On n'a aucune couverture par cycle qui utilise tous les arcs de $Y$.

  On a, pour une clause $c$, que 
  \[
    \sum_{\substack{\text{$C$ couverture par}\\ \text{cycles de $G_c$}}} \mathrm{poids}(C) = 4^3 \: \# \text{Assignations satisfaisant $c$}
  .\] 


  Au final, on obtient un graphe $G$ tel que
  \[
    \sum_{\substack{\text{$C$ couverture par}\\ \text{cycles de $G$}}} \mathrm{poids}(C) = 4^m \: \underbrace{\# \text{Assignations satisfaisant $\phi$}}_{\pbSharpSat(\phi)}
  ,\]
  où $m$ est le nombre d'occurrences de littéraux dans $\phi$.
  C'est ce facteur~$4^m$ qui fait que la réduction n'est pas parcimonieuse.


  \begin{rmk}
    Étant donnée une formule booléenne $\phi$, on construit une matrice  $A$ telle que $\perm(A) = 4^m  s$ $(*)$ où $s = \pbSharpSat(\phi)$.
    On a  $\perm(A) \mod 3 = s \mod 3$ ce qui est au moins aussi dur que $\np$ à calculer.

    Attention ! La relation $(*)$ ne dit rien sur la complexité pour calculer le permanent modulo $2$.
    En effet, on a \[
    \perm (X) \mod 2 = \det (X) \mod 2
    ,\] 
    ce qui est calculable en temps polynomial.
  \end{rmk}

  \section{Comptage modulaire.}

  \begin{thm}[Toda]
    On a $\ph \subseteq \oracle{\p}{\sharpp}$.
  \end{thm}

  \begin{defn}
    On définit $\intro*\modp$ par : un langage $A$ est dans $\modp$ ss'il existe un problème $B \in \p$ et un polynôme $p$ tels que \[
    x \in A \iff \# \{y \in \Sigma^{\le p(n)}  \mid \code{x,y} \in A\} \text{ est impair}
    .\]
  \end{defn}

  \begin{thm}
    On a $\np \subseteq \oracle{\rp}{\modp}$.\footnote{La classe $\rp$ a été vue en TD.}
    \label{thm:mod-thm1}
  \end{thm}

  \begin{rmk}
    La même preuve montre que le comptage modulo $3$ est $\np$-dur également.
  \end{rmk}

  \begin{lem}[d'isolation]
    On associe à chaque élément $x_i$ d'un ensemble $S = \{x_1, \ldots, x_n\} $ un poids entier $w_i \ge 1$.
    Soit $\mathcal{F}$ une famille de parties de $S$.
    Pour $T \in \mathcal{F}$, le poids $w(T)$ est défini par  $\sum_{i \text{ tq } x_i \in T} w_i$.
    Si les $w_i$ sont choisis indépendamment et suivant la distribution uniforme sur $\{1, \ldots, 2n\}$ alors, avec probabilité au moins $1/2$, il existe dans $\mathcal{F}$ un unique ensemble de poids minimaux.
  \end{lem}

  \begin{prv}
    On dit que $x_i \in S$ est "\textit{mauvais}" s'il appartient à un ensemble de $\mathcal{F}$ de poids minimum mais pas à tous.
    Il existe un unique élément de $\mathcal{F}$ de poids minimum ssi aucun $x_i$ n'est mauvais.

    On va montrer que, pour tout $i$, 
     \[
       \Pr[\text{$x_i$ est mauvais}] \le \frac{1}{2n}
    .\] 
    En effet, par borne de l'union, on peut conclure :
    \[
      \Pr[\exists i, \text{ $x_i$ est mauvais}] \le \frac{n}{2n} = \frac{1}{2}
    .\]
    On va montrer que, quels que soient les choix faits parmi les $w_j$ avec $j \neq i$, si on tire $w_i$ suivant une distribution uniforme sur $\{1, \ldots, 2n\}$ alors $\Pr[x_i \text{ est mauvais}] \le  1/2n$.

    Soit $w^\star(\bar{x}_i)$ le poids minimum d'un $T \in \mathcal{F}$ qui ne contient pas $x_i$, et posons
    \[
      w^\star(x_i) := \min_{T \text{ tq } x_i \in T} \sum_{\substack{j \neq i\\ x_j \in T}} w_j
    .\]
    On distingue deux cas.
    \begin{itemize}
      \item \textit{Cas 1} : $w^\star(x_i) \ge w^\star(\bar{x}_i)$.
        Alors on a $\Pr[x_i \text{ est mauvais}] = 0$ car $x_i$ n'appartient à aucun ensemble de poids minimaux.
      \item \textit{Cas 2} : $w^\star(\bar{x}_i) - w^\star(x_i) \ge 1$. On a que $x_i$ est mauvais ssi $w_i = w^\star(\bar{x}_i) - w^\star(x_i)$ 
        \begin{itemize}
          \item Si $w_i < w^\star(\bar{x}_i) - w^\star(x_i)$ alors $x_i$ est dans tous les ensembles de poids minimaux ;
          \item Si $w_i > w^\star(\bar{x}_i) - w^\star(x_i)$ alors $x_i$ n'est dans aucun ensemble de poids minimal.
        \end{itemize}
        Ainsi, on a que $w_i = w^\star(\bar{x}_i) - w^\star(x_i)$ est vrai avec probabilité $1/2n$ ou $0$.
    \end{itemize}
  \end{prv}

  \hspace{12pt}

  \begin{prv}[du théorème~\ref{thm:mod-thm1}]
    Il suffit de montrer que $\pbSat \in \oracle{\rp}{\modp}$.
    On considère l'algorithme suivant.

    \begin{algorithmic}[1]
      \Require une formule booléenne $F(x_1, \ldots, x_n)$
      \State On choisit, pour chaque $x_i$, un poids $w_i \in \{1, \ldots, 2n\}$.
      \State On choisit $u \in \{0, \ldots, 2n\}$.
      \State On note $w(\tau) = \sum_{i \text{ tq } \tau(x_i) = \texttt{1}} w_i$ pour une assignation $\tau$.
      \BeginBox[draw=nicered]
      \If{$\# \{\text{assignation $\tau$}  \mid \tau(F) = \texttt{1} \text{ et } w(\tau) \le u\} $ est impair}
      \State \textbf{Accepter}
      \Else
      \State \textbf{Rejeter}
      \EndIf
      \EndBox
    \end{algorithmic}
    La partie rouge utilise l'oracle pour $\modp$, car on peut vérifier en temps polynomial que $\tau(F) = \texttt{1}$ et que $w(\tau) \le u$.

    Analysons cet algorithme.
    \begin{enumerate}
      \item Si $F \not\in \pbSat$ alors $F$ est toujours refusée.
      \item Si  $F \in \pbSat$, on applique le lemme d'isolation à l'ensemble~$\mathcal{F}$ des $T \subseteq \{x_1, \ldots, x_n\}$ tels que l'assignation $\tau$ définie par~$\tau(x_i) = 1$ ssi $x_i \in T$ satisfait $F$.
        Avec probabilité au moins  $1 / 2$, on a une unique assignation $\tau$ satisfaisant $F$ de poids minimum $w^\star$.
        Si on prend $u = w^\star$, alors 
         \[
        \# \{\tau  \mid \tau(F) = \texttt{1} \text{ et } w(\tau) \le u\}  = 1
        ,\] 
        cela se produit avec probabilité au moins $1 / (2n^2 + 1)$.
        On a que $F$ est acceptée avec probabilité au moins $1 / 2 (2n^2 + 1)$.
        On peut répéter pour que la probabilité de succès soit au moins $1 / 2$.
    \end{enumerate}
  \end{prv}

  Le théorème~\ref{thm:mod-thm1} montre que $\np \subseteq \oracle{\rp}{\modp}$ en faisant plusieurs appels à l'oracle. Le théorème~\ref{thm:mod-thm2} améliore ça avec un seul appel à l'oracle.

  \begin{thm}
    On a $\np \subseteq \bp\modp$.
    \label{thm:mod-thm2}
  \end{thm}

  \begin{defn}
    Soient $A, R \subseteq \Sigma^\star$.
    On dit que $B$ est dans $\intro*\bp A$ s'il existe un polynôme $p$ tel que 
    \begin{enumerate}
      \item si $x \in A$ alors, $\# \{y \in \{\texttt{0},\texttt{1}\}^{p(n)}  \mid \code{x,y} \in \} \ge (3/4)\cdot 2^{p(n)}$ ;
      \item si $x \not\in A$ alors, $\# \{y \in \{\texttt{0},\texttt{1}\}^{p(n)}  \mid \code{x,y} \not\in \} \ge (3/4)\cdot  2^{p(n)}$ ;
    \end{enumerate}
    en notant $n = |x|$.
    
    C'est une réduction probabiliste de $B$ à $A$.

    Soit $\mathcal{C}$ une classe de complexité.
    On dit que $B \in \bp{\mathcal{C}}$ s'il existe $A \in \mathcal{C}$ tel que $B \in \bp A$, autrement dit, 
    \[
      \bp{\mathcal{C}} := \bigcup_{A \in \mathcal{C}}  \bp A
    .\]
  \end{defn}

  \begin{rmk}
    Les principales étapes de la preuve du théorème de Toda sont les deux inclusions
    \[
      \phSigma k \subseteq \bp\modp \subseteq \oracle \p \sharpp
    .\]
    La deuxième inclusion demande de \textit{supprimer l'aléatoire}.
  \end{rmk}

  \begin{prop}
    On a $\phSigma k \subseteq \bp\modp$.
  \end{prop}
  \begin{prv}[idée de la preuve]
    Par récurrence sur $k$, on montre que $\phSigma k \subseteq \bp{{\oplus \phPi {k-1}}}$ (pour le cas $k = 2$, c'est le théorème~\ref{thm:mod-thm2}).
    On a que $\phPi{k-1} \subseteq \bp \modp$, et $\bp\modp$ est clos par complément.
    En effet, on a donc
    \[
      \phSigma k \subseteq \bp {{\oplus \phPi {k-1}}} \subseteq \bp{{\oplus \bp{\modp}}} = \bp\modp
    .\]
  \end{prv}

  \begin{lem}
    On a que $A \in \modp$ ss'il existe un polynôme $p$ et un problème $B \in \p$ tel que, pour tout $x \in \Sigma^\star$,
    \[
      x \in A \iff \# \{y \in \{\texttt{0},\texttt{1}\}^{p(|x|)}  \mid \code{x, y} \in B\}  \text{ est pair}
    .\]
    Ainsi, $\modp$ est clos par complément.
  \end{lem}
  \begin{prv}
    Il existe $C \in \p$ et un polynôme $q$ tels que, pour toute entrée~$x \in \Sigma^\star$,
    \[
      x \in A \iff \# \{y \in \{\texttt{0},\texttt{1}\}^{q(|x|)}  \mid \code{x, y} \in C\}  \text{ est impair}
    .\]
    On définit $B$ par :
    \[
      B := \mleft\{\, \code{x, y} \;\middle|\;
      \begin{array}{c}
        y \neq \texttt{0}^{q(|x|)} \wedge \code{x,y} \in C\\
        \text{\textit{\textbf{ou}}}\\
        y = \texttt{0}^{q(|x|)} \wedge \code{x, y} \not\in C
      \end{array}\,\mright\} 
    .\]
    "On inverse la réponse sur $y = \texttt{0}^{q(|x|)}$."
  \end{prv}

  \begin{lem}
    Soit $A \in \modp$. L'ensemble
    \[
      A' := \Big\{\, \code{x_1, \ldots, x_k} \;\Big|\; \bigvee_{i=1}^k (x_i \in A)\,\Big\} 
    \] est dans $\modp$.
  \end{lem}
  \begin{prv}
    Par le lemme précédent, il existe $p$ un polynôme et $B \in \p$ tel que,
    \[
      x \in A \iff \# \{y \in \{\texttt{0},\texttt{1}\}^{p(|x|)}  \mid \code{x, y} \in B\}  \text{ est pair}
    .\] 
    On pose $B'$ l'ensemble
    \[
      B' := \Big\{\,\code{x_1, \ldots, x_k, y_1, \ldots, y_k} \;\Big|\; \bigwedge_{i=1}^k (\code{x_i, y_i} \in B \,\Big\} 
    .\]
    On a 
    \begin{align*}
      &\code{x_1, \ldots, x_k} \in A' \\
      \iff& \bigvee_{i=1}^k \# \{y_i \in \{\texttt{0},\texttt{1}\}^{p(|x|}  \mid \code{x_i, y_i} \in B \} \text{ pair}\\
      \iff& \# \{ \code{y_1, \ldots, y_n}  \mid  \code{x_1, \ldots, x_k, y_1, \ldots, y_k} \in B\} \text{ pair} \\
    .\end{align*}
  \end{prv}

  \hspace{12pt}

  \begin{prv}[du théorème~\ref{thm:mod-thm2}]
    On montre que $\pbSat \in \bp\modp$.
    Soit $F(x_1, \ldots, x_n)$ une formule booléenne et soit $m = 2 (2n^2 + 1)$.
    On accepte $F$ ssi
    \[
    \code{F, w^{(1)}, u_1, w^{(2)}, u_2, \ldots, w^{(2m)}, u_{2m}}
    \]
    (où les $w^{(i)}$ et les $u_i$ sont tirés au sort)
    satisfait la propriété
    \[
    \bigvee_{j=1}^{2m} \# \{\tau \in \{\texttt{0},\texttt{1}\}^n  \mid \tau(F) = \texttt{1} \text{ et } w^{(j)}(\tau) \le u_j\} \text{ est impair}
    ,\] 
    et ce dernier est dans $\modp$ par le lemme précédent.

    Analysons cet algorithme.
    \begin{itemize}
      \item Si $F \not\in \pbSat$, alors on accepte $F$ avec probabilité $0$ car les $2m$ expériences échouent.
      \item Si $F \in \pbSat$, alors on a
        \[
          \Pr[F \text{ rejetée}] \le \left( 1 - \frac{1}{m} \right)^{2m} \le \frac{1}{e^2}
        ,\] 
        d'où $\Pr[F \text{ acceptée}] \ge 3 / 4$.
    \end{itemize}
  \end{prv}
\end{document}
